\documentclass{book}
\usepackage[utf8]{inputenc}
\usepackage[T1]{fontenc}
\usepackage{geometry}
\usepackage{graphicx}
\usepackage{float}
\usepackage{hyperref}
\usepackage[english,indonesian]{babel}

\title{Optimasi Pemrosesan Template LaTeX Menggunakan Pendekatan Berbasis AI}
\author{Academic Support Agent}
\date{\today}

\begin{document}
\maketitle

\begin{abstract}

\end{abstract}

\chapter{PENDAHULUAN}


\chapter{TINJAUAN PUSTAKA}


% Original content (commented out):
% 
%%%%%%%%%%%%%%%%%%%%%%%%%%%레이텍스 기본 설정%%%%%%%%%%%%%%%%%%%%%%%%%%%%%%%%%
\documentclass[12pt, twoside]{book} %문서의 형식을 report로 지정하고 글자포인트를 지정한다. %book, report, slide, letter %towside는 1장씩 보는 문서가 아닌 2장 씩 보는 문서로 만들게 해준다. 
\usepackage[utf8]{inputenc}
\usepackage[a4paper, width=170mm,top=15mm,bottom=15mm]{geometry}
%a4paper는 미국 문서 양식이다. twoside로 설정했시때문에 매 장마다 문서가 좌우 살짝씩 다르다. width는 문서의 넓이를 조정하고, top과 bottom은 각 세션의 마진을 정한다.
%%%%%%%%%%%%%%%%%%%%%%%%%%%%%%%%%%%%%%%%%%%%%%%%%%%%%%%%%%%%



%%%%%%%%%%%%%%%%%%%%%%%%%%여러가지 패키지들%%%%%%%%%%%%%%%%%%%%%%%%%%%%%%%%%%

\usepackage{graphicx}%그래픽(사진)을 첨부할 수 있도록 해주는 패키지이다.
\graphicspath{{images/}} %이미지가 있는 경로를 지정해준다.
\usepackage{float}%그림을 페이지에 고정하도록 한다. [H]
\usepackage{caption}
\usepackage{subcaption} %\usepackage{caption} 와 \usepackage{subcaption}는 세개의 그림을 하나의 라인에 같이 나란히 배열할 수 있도록 해주는 패키지이다.
\usepackage{amsmath}%수학 기호 작성을 쉽게해주는 패키지이다.
\usepackage{linguex}%Interlinear Glossing을 할 수 있게 해주는 패키지이다.
\usepackage{kotex}%한국어를 사용할 수 있게 해준다.
\usepackage{indentfirst}%들여쓰기를 할 수 있게 해준다.
%%%%%%%%%%%%%%%%%%%%%%%%%%%%%%%%%%%%%%%%%%%%%%%%%%%%%%%%%%%%


%%%%%%%%%%%%%%%%%%%%%%%%%%%%%폰트 스타일%%%%%%%%%%%%%%%%%%%%%%%%%%%%%%%
\usepackage[sfdefault]{roboto}  %% Option 'sfdefault' only if the base font of the document is to be sans serif
% \usepackage{helvet}
\usepackage[T1]{fontenc}
\usepackage{libertine}
% 볼드 사이즈 지정
\newcommand{\textBF}[1]{%
    \pdfliteral direct {2 Tr 0.35 w} %the second factor is the boldness
     #1%
    \pdfliteral direct {0 Tr 0 w}%
}
%%%%%%%%%%%%%%%%%%%%%%%%%%%%%%%%%%%%%%%%%%%%%%%%%%%%%%%%%%%%


%original
%%%%%%%%%%%%%%%%%%%%%%%%%%%%하이퍼링크 생성%%%%%%%%%%%%%%%%%%%%%%%%%%%%%%%
\usepackage[dvipsnames]{xcolor}
\definecolor{DDbule}{RGB}{45, 72, 207}
%하이퍼링크 색상을 더 추가하게 해준다.

\usepackage{hyperref}
\hypersetup{
    colorlinks=true,
    linkcolor=DDbule,
    citecolor=DDbule,
    urlcolor=DDbule,
    linktoc=page %목록에서 페이지 번호에 링크가 동작하도록 한다.
}
% linktoc | =none,section,page,all | defines which part of an entry in the table of contents is made into a link
%%%%%%%%%%%%%%%%%%%%%%%%%%%%%%%%%%%%%%%%%%%%%%%%%%%%%%%%%%%







%%%%%%%%%%%%%%%%%%%%%%%%%%각 챕터 상단과 하단 스타일 지정%%%%%%%%%%%%%%%%%%%%%%%%%%%%%%%%%%
\usepackage{fancyhdr}
\pagestyle{fancy}
\newcommand{\changefont}{%
    \fontsize{10}{12}\selectfont
}
%상단 챔터 및 세션 정보의 폰트를 지정한다.
%https://www.overleaf.com/learn/latex/headers_and_footers
\fancyhf{}
% \fancyhead[LO]{\changefont \rightmark} %section
% \fancyhead[RE]{\changefont \leftmark} % chapter
%revised: 페이지 번호 위치
\fancyhead[LE]{\changefont \rightmark} %section
\fancyhead[RO]{\changefont \leftmark} % chapter


%fancyhdr패키지를 사용하면 각 챕터마다 상단에 스타일이 들어간다.챕터이름, 세션이름

%fancyhdr는 세부적으로 스타일을 조정할 수있다.
%\fancyhead{}
%\fancyhead[LO,RE]{Thesis Title}
%\fancyhead[LO,RE]{}

%\fancyhead[LO,RE]{Thesis Title}는 헤더부분에 지정한 텍스트를 왼쪽으로 쏠린 장에서 왼쪽에 LO, 오른쪽으로 쏠린장에서 오른쪽에 표기한다. RE

%\fancyfoot[LO,RE]{\thepage}
%revised: 페이지 번호 위치
\fancyfoot[LE,RO]{\thepage}

%\fancyfoot[LE,RO]{\thepage}는 하단부분에 해당 페이짖 번호를 왼쪽으로 쏠린 장에서 오른쪽에 LE, 오른쪽으로 쏠린장에서 왼쪽에 표기한다. RO
%\fancyfoot[CE,CO]{\thepage}
%\fancyfoot[CE,CO]{\thepage}는 페이지 번호를 가운데에 표기한다.
%\thechapter는 해당 챕터의 이름을 표기한다.
%\leftmark는 챕터와 챕터 제목을 표기한다.
%\rightmark는 해당 페이지의 세션을 지정한다.

%\fancyfoot{}
%\fancyfoot[LE,RO]{\thepage}
%\fancyfoot[LO,CE]{Chapter \thechapter}
%\fancyfoot[CO,RE]{Seongmin Mun}
%\renewcommand{\headrulewidth}{0.4pt}
%\renewcommand{\footrulewidth}{0.4pt}
%\renewcommand{\headrulewidth}{0.4pt}로 헤더에 줄을 추가하고 사이즈를 조정한다.
%\renewcommand{\footrulewidth}{0.4pt}로 하단에 줄을 추가하고 사이즈를 조정한다.
%%%%%%%%%%%%%%%%%%%%%%%%%%%%%%%%%%%%%%%%%%%%%%%%%%%%%%%%%%%%






%%%%%%%%%%%%%%%%%%%%%%%%%%%%%챕터 스타일%%%%%%%%%%%%%%%%%%%%%%%%%%%%%%%
\usepackage[Lenny]{fncychap} %각 챕터의 스타일을 지정한다.
%종류는 Sonny, Lenny, Glenn, Conny, Rejne and Bjarne.이 있다.
%http://tug.ctan.org/macros/latex/contrib/fncychap/fncychap.pdf 참고
%챕터 제목, 번호 등등 사용자지정하는 법
%\ChNameVar{\Large\fontfamily{put}\selectfont\color{blue}}
%\ChNumVar{\Huge\fontfamily{pzc}\selectfont\color{red}}
%\ChTitleVar{\Large\fontfamily{phv}\selectfont\scshape\color{green}}
\ChTitleVar{\LARGE\bfseries} %글씨크기, 볼드체
%\Huge
%\huge
%\LARGE
%\Large
%\large
%\normalsize (default)
%\small
%\footnotesize
%\scriptsize
%\tiny
%%%%%%%%%%%%%%%%%%%%%%%%%%%%%%%%%%%%%%%%%%%%%%%%%%%%%%%%%%%%





%%%%%%%%%%%%%%%%%%%%%%%%%%%문장 사이 간격조정%%%%%%%%%%%%%%%%%%%%%%%%%%%%%%%%%
\usepackage{setspace}
\doublespacing
%문단을 두줄 문단으로 바꾸어 준다.
% Single line spacing, alternatives: onehalfspacing or doublespacing
%%%%%%%%%%%%%%%%%%%%%%%%%%%%%%%%%%%%%%%%%%%%%%%%%%%%%%%%%%%%


%%%%%%%%%%%%%%%%%%%%%%%%%%%컴파일 속도%%%%%%%%%%%%%%%%%%%%%%%%%%%%%%%%%
\usepackage{chemformula}%기본 설정 컴파일 속도보다 빠르게 해준다.
\let\ce\ch
%%%%%%%%%%%%%%%%%%%%%%%%%%%%%%%%%%%%%%%%%%%%%%%%%%%%%%%%%%%%




%%%%%%%%%%%%%%%%%%%%%%%%%%%%%%레퍼런스%%%%%%%%%%%%%%%%%%%%%%
\usepackage[sort,comma,authoryear,round]{natbib}%참고문헌의 스타일을 더 다양하게 만든다.
%sort를 사용하여 저자 목록을 자동으로 정렬해준다.
\usepackage[nottoc,notlof,notlot]{tocbibind} 
\renewcommand\bibname{References}
%레퍼런스 이름을 bibliographic에서 reference로 바꾼다.
%%%%%%%%%%%%%%%%%%%%%%%%%%%%%%%%%%%%%%%%%%%%%%%%%%%%%%%%%%%%






%%%%%%%%%%%%%%%%%%%%%%%%%%%%%%이전 타이틀이름%%%%%%%%%%%%%
%\title{
%    {Thesis Title}\\
%    {\large Institution Name}\\
%    {\includegraphics{BERT.png}}
%}

%\title{Thesis Title}
%\author{Seongmin Mun}
%\date{February 2021}
%%%%%%%%%%%%%%%%%%%%%%%%%%%%%%%%%%%%%%%%%%%%%%%%%%%%%%%%%%%%







%%%%%%%%%%%%%%%%%%%%%%%%%목록깊이%%%%%%%%%%%%%%%%%%%%%%%%%%%%%%%%%%%
%%%목록의 깊이를 3번까지 본다.
\setcounter{tocdepth}{3}
%%%%%%%%%%%%%%%%%%%%%%%%%%%%%%%%%%%%%%%%%%%%%%%%%%%%%%%%%%%%





%%%%%%%%%%%%%%%%%%%%%%%%%테이블 여백을 지정%%%%%%%%%%%%%%%%%%%%%%%%%%%%%%%%%%%
\setlength{\tabcolsep}{10pt}
\renewcommand{\arraystretch}{1.3}
\usepackage{ctable} %테이블 중 굵기를 조정하게 해준다.
%\specialrule{.05em}{1em}{0em} 
\usepackage{longtable} %테이블 다음 페이지까지 연장
\usepackage{tabularx, blindtext}
%%%%%%%%%%%%%%%%%%%%%%%%%%%%%%%%%%%%%%%%%%%%%%%%%%%%%%%%%%%%





%%%%%%%%%%%%%%%%%%%%%%%%%테이블 정렬 및 넓이 조정%%%%%%%%%%%%%%%%%%%%%%%%%%%%%%%%%%%
\usepackage{array}
\newcommand{\PreserveBackslash}[1]{\let\temp=\\#1\let\\=\temp}
\newcolumntype{C}[1]{>{\PreserveBackslash\centering}p{#1}}
\newcolumntype{R}[1]{>{\PreserveBackslash\raggedleft}p{#1}}
\newcolumntype{L}[1]{>{\PreserveBackslash\raggedright}p{#1}}
\usepackage{multirow} %행 병합
%%%%%%%%%%%%%%%%%%%%%%%%%%%%%%%%%%%%%%%%%%%%%%%%%%%%%%%%%%%%%%%%%%%%%%%%%%%%%%%








%original
%%%%%%%%%%%%%%%%%%%%%%%%%%코드 스타일%%%%%%%%%%%%%%%%%%%%%%%%%%%%%%%%%%%
\usepackage{color}
\definecolor{lightgray}{rgb}{0.95, 0.95, 0.95}
\definecolor{darkgray}{rgb}{0.4, 0.4, 0.4}
%\definecolor{purple}{rgb}{0.65, 0.12, 0.82}
\definecolor{editorGray}{rgb}{0.95, 0.95, 0.95}
\definecolor{editorOcher}{rgb}{1, 0.5, 0} % #FF7F00 -> rgb(239, 169, 0)
\definecolor{editorGreen}{rgb}{0, 0.5, 0} % #007C00 -> rgb(0, 124, 0)
\definecolor{orange}{rgb}{1,0.45,0.13}		
\definecolor{olive}{rgb}{0.17,0.59,0.20}
\definecolor{brown}{rgb}{0.69,0.31,0.31}
\definecolor{purple}{rgb}{0.38,0.18,0.81}
\definecolor{lightblue}{rgb}{0.1,0.57,0.7}
\definecolor{lightred}{rgb}{1,0.4,0.5}
\usepackage{upquote}
\usepackage{listings}
% CSS
\lstdefinelanguage{CSS}{
  keywords={color,background-image:,margin,padding,font,weight,display,position,top,left,right,bottom,list,style,border,size,white,space,min,width, transition:, transform:, transition-property, transition-duration, transition-timing-function},	
  sensitive=true,
  morecomment=[l]{//},
  morecomment=[s]{/*}{*/},
  morestring=[b]',
  morestring=[b]",
  alsoletter={:},
  alsodigit={-}
}

% JavaScript
\lstdefinelanguage{JavaScript}{
  morekeywords={typeof, new, true, false, catch, function, return, null, catch, switch, var, if, in, while, do, else, case, break},
  morecomment=[s]{/*}{*/},
  morecomment=[l]//,
  morestring=[b]",
  morestring=[b]'
}

\lstdefinelanguage{HTML5}{
  language=html,
  sensitive=true,	
  alsoletter={<>=-},	
  morecomment=[s]{<!-}{-->},
  tag=[s],
  otherkeywords={
  % General
  >,
  % Standard tags
	<!DOCTYPE,
  </html, <html, <head, <title, </title, <style, </style, <link, </head, <meta, />,
	% body
	</body, <body,
	% Divs
	</div, <div, </div>, 
	% Paragraphs
	</p, <p, </p>,
	% scripts
	</script, <script,
  % More tags...
  <canvas, /canvas>, <svg, <rect, <animateTransform, </rect>, </svg>, <video, <source, <iframe, </iframe>, </video>, <image, </image>, <header, </header, <article, </article
  },
  ndkeywords={
  % General
  =,
  % HTML attributes
  charset=, src=, id=, width=, height=, style=, type=, rel=, href=,
  % SVG attributes
  fill=, attributeName=, begin=, dur=, from=, to=, poster=, controls=, x=, y=, repeatCount=, xlink:href=,
  % properties
  margin:, padding:, background-image:, border:, top:, left:, position:, width:, height:, margin-top:, margin-bottom:, font-size:, line-height:,
	% CSS3 properties
  transform:, -moz-transform:, -webkit-transform:,
  animation:, -webkit-animation:,
  transition:,  transition-duration:, transition-property:, transition-timing-function:,
  }
}

\lstdefinestyle{htmlcssjs} {%
  % General design
%  backgroundcolor=\color{editorGray},
  basicstyle={\footnotesize\ttfamily},   
  frame=b,
  % line-numbers
  xleftmargin={0.75cm},
  numbers=left,
  stepnumber=1,
  firstnumber=1,
  numberfirstline=true,	
  % Code design
  identifierstyle=\color{black},
  keywordstyle=\color{blue}\bfseries,
  ndkeywordstyle=\color{editorGreen}\bfseries,
  stringstyle=\color{editorOcher}\ttfamily,
  commentstyle=\color{brown}\ttfamily,
  % Code
  language=HTML5,
  alsolanguage=JavaScript,
  alsodigit={.:;},	
  tabsize=2,
  showtabs=false,
  showspaces=false,
  showstringspaces=false,
  extendedchars=true,
  breaklines=true,
  % German umlauts
  literate=%
  {Ö}{{\"O}}1
  {Ä}{{\"A}}1
  {Ü}{{\"U}}1
  {ß}{{\ss}}1
  {ü}{{\"u}}1
  {ä}{{\"a}}1
  {ö}{{\"o}}1
}
%
\lstdefinestyle{py} {%
language=python,
literate=%
*{0}{{{\color{lightred}0}}}1
{1}{{{\color{lightred}1}}}1
{2}{{{\color{lightred}2}}}1
{3}{{{\color{lightred}3}}}1
{4}{{{\color{lightred}4}}}1
{5}{{{\color{lightred}5}}}1
{6}{{{\color{lightred}6}}}1
{7}{{{\color{lightred}7}}}1
{8}{{{\color{lightred}8}}}1
{9}{{{\color{lightred}9}}}1,
basicstyle=\footnotesize\ttfamily, % Standardschrift
numbers=left,               % Ort der Zeilennummern
%numberstyle=\tiny,          % Stil der Zeilennummern
%stepnumber=2,               % Abstand zwischen den Zeilennummern
numbersep=5pt,              % Abstand der Nummern zum Text
tabsize=4,                  % Groesse von Tabs
extendedchars=true,         %
breaklines=true,            % Zeilen werden Umgebrochen
keywordstyle=\color{blue}\bfseries,
frame=b,
commentstyle=\color{brown}\itshape,
stringstyle=\color{editorOcher}\ttfamily, % Farbe der String
showspaces=false,           % Leerzeichen anzeigen ?
showtabs=false,             % Tabs anzeigen ?
xleftmargin=17pt,
framexleftmargin=17pt,
framexrightmargin=5pt,
framexbottommargin=4pt,
%backgroundcolor=\color{lightgray},
showstringspaces=false,      % Leerzeichen in Strings anzeigen ?
}%
%%%%%%%%%%%%%%%%%%%%%%%%%%%%%%%%%%%%%%%%%%%%%%%%%%%%%%%%%%%%%%%%%%%%%%%%%%%%%%%%











%%%%%%%%%%%%%%%%%%%%%%%%%%%%%%%%%%%%%%%%%%문서의시작 %%%%%%%%%%%%%%%%%%%%%%%%%%%%

\begin{document}


%%첫번째 챕터가 시작되기 전 파트에는 챕터 상하단 스타일을 제거한다.
\frontmatter
\pagestyle{empty}

% \newgeometry{width=210mm,top=0mm,bottom=0mm}

% %타이틀 파일을 따로 열어서 생성한다.
% \input{coverpage}

% \newgeometry{width=170mm,top=15mm,bottom=15mm}

%타이틀 파일을 따로 열어서 생성한다.
\input{titlepage}

%타이틀 스타일과 다른 스타일로 다른 세션을 진행한다.
\newgeometry{width=140mm,top=40mm,bottom=40mm}

%\maketitle

%abstract를 파일로 사용하려면 아래를 사용
%\input{abstract}

\chapter*{Acknowledgements}
\noindent There are many who helped me along the way on this journey. 

\chapter*{Abstract}
\noindent This dissertation reports computational accounts of resolving word-level polysemy in a lesser-studied language—Korean. 
\newline
\indent \textBF{Keywords:} polysemy, natural language processing, classification, word embedding models, data visualization, Korean


\tableofcontents
%\tableofcontents는 문서 전체에서 사용된 챕터를 목록으로 만들어 준다.


\listoftables
%\listoftables는 문서 전체에서 사용된 표를 목록으로 만들어 준다.


\listoffigures
%\listoffigures는 문서 전체에서 사용된 그림을 목록으로 만들어 준다.

% \lstlistoflistings
% %\lstlistoflistings는 문서 전체에서 사용된 코드를 목록으로 만들어 준다.



\chapter*{List of abbreviations}
\noindent I follow the Leipzig glossing rules\footnote{Available at: \href{https://www.eva.mpg.de/lingua/pdf/Glossing-Rules.pdf}{https://www.eva.mpg.de/lingua/pdf/Glossing-Rules.pdf}}.
\\
\begin{table}[htp!]
    \centering
    \begin{tabular}{ll} %l 외쪽 c 가운데 r오른쪽
        \specialrule{.08em}{0.2em}{0em} 
        \textbf{Abbreviation} & \textbf{Label} \\ %\\는 다음 행으로 넘어가게 해준다.
        \specialrule{.08em}{0.2em}{0em} 
        ACC & Accusative \\
        \hline
        AGT & Agent \\
        \hline
    \end{tabular}
\end{table}





 
%다시 스타일을 적용한다.
\mainmatter
\pagestyle{fancy}

%챕터만 가운데로 표기되던 페이지 번호를 다른 형식들과 일치시켜 바꿔준다.
% \fancypagestyle{plain}{%
%   \fancyhead{}%
%   \renewcommand*{\headrule}{}%
%   \fancyfoot{}%
%   \fancyfoot[LO,RE]{\thepage}%
% }

%revised: 페이지 번호 위치
\fancypagestyle{plain}{%
  \fancyhead{}%
  \renewcommand*{\headrule}{}%
  \fancyfoot{}%
  \fancyfoot[LE,RO]{\thepage}%
}


\chapter{Introduction}\label{chap:chap1}
\input{chapters/chapter/chapter01}
\clearpage{\pagestyle{empty}\cleardoublepage}%챕터 뒤에 공백에 스타일 적용을 없앤다





\renewcommand{\theHchapter}{A\arabic{chapter}}
\appendix
\chapter{Code for the word-level embedding models}\label{append:append2}
\input{chapters/appendix/appendix01.tex}
\clearpage{\pagestyle{empty}\cleardoublepage}



\bibliographystyle{plainnat}  %참고문헌 목록의 스타일을 지정한다. 이것보다 많은 스타일이 있다. https://www.overleaf.com/learn/latex/bibtex_bibliography_styles
\bibliography{reference}

\end{document}





%%%%   1

%2개의 폴더를 만든다.images, chapters
%문서의 형식을 report로 지정하고 글자포인트를 지정한다.
%\usepackage{graphicx}%그래픽(사진)을 첨부할 수 있도록 해주는 패키지이다.
%\graphicspath{{images/}} %이미지가 있는 경로를 지정해준다.
%title을 재설정해준다.
%abstract를 넣어준다.
%Dedication를 넣어준다.
%Declaration를 넣어준다.
%Acknowledgements를 넣어준다.
%5개의 챕터를 만들어준다.
%각각 연결해준다.
%abstract를 넣어준다.
%chapter에 각 챕터의 제목을 넣고 input으로 파일을 연결




%%%%   2

%towside는 1장씩 보는 문서가 아닌 2장 씩 보는 문서로 만들게 해준다. 
%a4paper는 미국 문서 양식이다. twoside로 설정했시때문에 매 장마다 문서가 좌우 살짝씩 다르다. width는 문서의 넓이를 조정하고, top과 bottom은 각 세션의 마진을 정한다.
%fancyhdr패키지를 사용하면 각 챕터마다 상단에 스타일이 들어간다.챕터이름, 세션이름
%fancyhdr는 세부적으로 스타일을 조정할 수있다.
%\fancyhead[LO,RE]{Thesis Title}는 헤더부분에 지정한 텍스트를 왼쪽으로 쏠린 장에서 왼쪽에 LO, 오른쪽으로 쏠린장에서 오른쪽에 표기한다. RE
%\fancyfoot[LE,RO]{\thepage}는 하단부분에 해당 페이짖 번호를 왼쪽으로 쏠린 장에서 오른쪽에 LE, 오른쪽으로 쏠린장에서 왼쪽에 표기한다. RO
%\fancyfoot[CE,CO]{\thepage}는 페이지 번호를 가운데에 표기한다.
%\renewcommand{\headrulewidth}{0.4pt}로 헤더에 줄을 추가하고 사이즈를 조정한다.
%\renewcommand{\footrulewidth}{0.4pt}로 하단에 줄을 추가하고 사이즈를 조정한다.

%############팁##########만약에 chapter02이후 모든 챕터에서 스타일을 취소하고 싶으면 chapter02에서 \pagestyle{empty}를 추가한다.
%이것 외에 plain은 번호만 들어가게 해주고 다시 \pagestyle{fancy}를 입력하면 원상복구 된다.





%   3
%챕터 2에 아래와 같이 작성한다.
%qui officia deserunt mollit anim id est laborum. \ref{fig:BERT}
%\begin{figure}
%    \centering
%    \includegraphics[scale=0.5]{images/BERT.png}
%    \caption{BERT caption}
%    \label{fig:BERT}
%\end{figure}
%각 챕터에 그림을 삽입하는 방법이다.

%다음은 세개의 그림을 하나의 라인에 같이 나란히 배열하는 방법이다.
%\usepackage{caption} 와 \usepackage{subcaption}는 세개의 그림을 하나의 라인에 같이 나란히 배열할 수 있도록 해주는 패키지이다.
%챕터 3에 가서 아래와 같이 작성한다.

%\begin{figure}
%    \centering
%    \begin{subfigure}[b]{0.3\textwidth} %아래쪽에 배열하고 문서 넓이의 3분의 1배이다.
%        \centering
%        \includegraphics[width=\textwidth]{images/impact_b.png}
%        \caption{impact caption}
%        \label{fig:impact}
%    \end{subfigure}
%    \hfill   %그림간 간격을 조정해준다.
%    \begin{subfigure}[b]{0.3\textwidth} %아래쪽에 배열하고 문서 넓이의 3분의 1배이다.
%        \centering
%        \includegraphics[width=\textwidth]{images/person_b.png}
%        \caption{person caption}
%        \label{fig:person}
%    \end{subfigure}
%    \hfill   %그림간 간격을 조정해준다.
%    \begin{subfigure}[b]{0.3\textwidth} %아래쪽에 배열하고 문서 넓이의 3분의 1배이다.
%        \centering
%        \includegraphics[width=\textwidth]{images/question_b.png}
%        \caption{question caption}
%        \label{fig:question}
%    \end{subfigure}
%    \caption{mark mark}
%    \label{fig:marks}
%\end{figure}


%\listoffigures는 문서 전체에서 사용된 그림을 목록으로 만들어 준다.

%chapter04에 테이블을 작성한 것을 참고
%\listoftables는 문서 전체에서 사용된 표를 목록으로 만들어 준다.





%    4

%\usepackage[round]{natbib}%참고문헌의 스타일을 더 다양하게 만든다.
%\bibliographystyle{plainnat}  %참고문헌 목록의 스타일을 지정한다. 이것보다 많은 스타일이 있다. https://www.overleaf.com/learn/latex/bibtex_bibliography_styles
%\bibliography{example}


%natbib패키지를 사용하면 더 효율적으로 참고문헌의 스타일을 바꿀 수 있다.
%\usepackage[round]{natbib}%참고문헌의 스타일을 더 다양하게 만든다.
%\citep를 사용할 수 있다.
%\bibliographystyle{plain}를 \bibliographystyle{plainnat}로 바꾸어 natbib스타일로 사용한다.




%    5
% title, abstract를 따로 만들기



\end{document}